\documentclass{article}
\usepackage[utf8]{inputenc}
\usepackage{amsmath}
\usepackage{booktabs}
\usepackage{url}

\begin{document}

  \section{Punto flotante de 8 bits}

    Para los números de punto flotante de 32 y 64 bits existe una
    definición exacta dada por el estándar IEEE 754. FP8, en cambio, no
    forma parte de ese estándar; la propuesta más difundida es un formato
    de intercambio con dos codificaciones, E4M3 y E5M2, 
    también recogido en la especificación OFP8
    del Open Compute Project. En la notación E$x$M$y$,
    $x$ es el número de bits del exponente e $y$ el de la mantisa, y hay
    además 1 bit de signo.

    \subsection{Formatos}

      Ambos formatos ocupan 8 bits y difieren en cómo los reparten. E4M3
      destina más bits a la mantisa y por tanto tiene mayor precisión; E5M2
      destina más bits al exponente y por tanto tiene mayor rango dinámico.
      Otra diferencia es el tratamiento de los valores especiales: E5M2
      sigue las convenciones de IEEE 754 (con $\pm\infty$ y varios NaN),
      mientras que E4M3 prescinde de los infinitos y reserva un único patrón
      de mantisa para NaN, lo que le permite ampliar su rango máximo.

      \begin{center}
        \begin{tabular}{lcc}
          \toprule
          & E4M3 & E5M2 \\
          \midrule
          Bits (signo/exponente/mantisa) & 1/4/3 & 1/5/2 \\
          Sesgo del exponente & 7 & 15 \\
          Infinitos & no & sí \\
          Menor positivo (subnormal) & $2^{-9}$ & $2^{-16}$ \\
          Mayor valor finito & 448 & 57344 \\
          Error relativo máximo de redondeo & $6.25\,\%$ & $12.5\,\%$ \\
          \bottomrule
        \end{tabular}
      \end{center}

    \subsection{Uso de esta representación por los desarrolladores de DeepSeek}

      En el informe técnico de DeepSeek-V3 \cite{deepseek2024v3}, sus
      desarrolladores describen un esquema de entrenamiento de
      \emph{precisión mixta} en FP8, que validaron por primera vez en un
      modelo de escala muy grande. La idea central es que las operaciones
      que concentran el cómputo se ejecutan en 8 bits, mientras que unas
      pocas operaciones sensibles conservan formatos más anchos (BF16 o
      FP32). A continuación se resume qué formato se usa en cada parte del
      cálculo y qué se obtiene con cada elección.

      \subsubsection{E4M3 en todas las multiplicaciones matriciales}

        En una capa lineal intervienen tres multiplicaciones matriciales
        (GEMM): la pasada hacia adelante (\texttt{Fprop}), la
        retropropagación respecto a la activación (\texttt{Dgrad}) y la
        retropropagación respecto a los pesos (\texttt{Wgrad}). DeepSeek
        ejecuta las tres en FP8, con entradas en E4M3 y salidas en BF16 o
        FP32 \cite{deepseek2024v3}. Los beneficios que reportan son:

        \begin{itemize}
          \item \emph{Velocidad:} en teoría, el cómputo se duplica
            respecto al que se obtendría con BF16.

          \item \emph{Precisión:} E4M3 tiene 3 bits de mantisa frente a los
            2 de E5M2, es decir, un error relativo máximo de redondeo de
            $6.25\,\%$ en lugar de $12.5\,\%$.

        \end{itemize}

      \subsubsection{Por qué no se usa E5M2}

        Los esquemas híbridos de trabajos previos emplean E4M3 en
        \texttt{Fprop} y E5M2 en \texttt{Dgrad} y \texttt{Wgrad}
        \cite{deepseek2024v3,micikevicius2022fp8}. La ventaja de E5M2 es
        su rango dinámico (hasta 57344, frente a 448 de E4M3), útil para
        valores con una distribución muy dispersa, a costa de una
        mantisa más corta. DeepSeek toma la decisión contraria, que el
        informe resume como priorizar la mantisa sobre el exponente: usa
        E4M3 en todos los tensores para ganar precisión. En el esquema
        descrito, E5M2 no se emplea en ninguna de las tres GEMM; aparece
        solo como la alternativa que se descarta.

        Los autores atribuyen la viabilidad de esta elección al escalado
        de grano fino, que se explica a continuación: al escalar grupos
        pequeños de elementos, los grupos comparten de hecho un exponente,
        lo que mitiga el rango limitado de E4M3.

      \subsubsection{Partes que se mantienen en mayor precisión}

        No todo pasa a 8 bits. El informe indica que se conserva BF16 o
        FP32 en el módulo de embeddings, la cabeza de salida, los módulos
        de \emph{gating} de la mezcla de expertos (MoE), las operaciones de
        normalización y los operadores de atención. Además, los pesos
        maestros y los gradientes de los pesos se guardan en FP32, y los
        dos momentos del optimizador AdamW en BF16 \cite{deepseek2024v3}.
        El beneficio buscado es la estabilidad numérica del entrenamiento;
        el costo en memoria de estos componentes se reduce repartiéndolos
        entre los procesos de paralelismo de datos.

      \subsubsection{Almacenamiento y comunicación}

        Los formatos de 8 bits también se usan fuera de las GEMM para
        ahorrar memoria y ancho de banda \cite{deepseek2024v3}:

        \begin{itemize}
          \item Las activaciones que necesita \texttt{Wgrad} se guardan
            en FP8, y las entradas de la función SwiGLU en la MoE también,
            recalculando su salida en la pasada hacia atrás.
          \item Como excepción, las entradas de la capa lineal posterior a
            la atención usan un formato propio E5M6 (12 bits), porque
            también intervienen en la retropropagación de la atención y son
            sensibles a la precisión.
          \item En la comunicación de la MoE, la activación previa a las
            proyecciones de expansión se cuantiza a FP8 antes del
            \emph{dispatch}, lo que reduce el tráfico entre nodos; la
            operación de \emph{combine}, en cambio, se mantiene en BF16
            para preservar la precisión.
        \end{itemize}

      \subsubsection{Evolución posterior: escalas UE8M0}

        En la versión DeepSeek-V3.1, la ficha del modelo indica que el
        entrenamiento usó el formato de escala UE8M0 FP8 en pesos y
        activaciones, por compatibilidad con los formatos de microescalado
        \cite{deepseek2025v31}. UE8M0 no codifica los datos sino los
        factores de escala: es un formato sin signo, con 8 bits de
        exponente y ninguno de mantisa, de modo que solo representa
        potencias de dos y cada escala ocupa un byte. Esto continúa una
        línea ya presente en V3, donde algunos factores de escala se
        restringían a potencias enteras de dos \cite{deepseek2024v3}.


      Como resumen, en el esquema descrito E4M3 aparece en todos los
      cálculos en 8 bits de las GEMM y E5M2 no se utiliza. La mayor
      precisión de E4M3 se compensa con el escalado por grupos y con la
      acumulación en FP32, y las partes más sensibles permanecen en
      formatos más anchos.

  \begin{thebibliography}{9}

    \bibitem{micikevicius2022fp8}
      P.~Micikevicius, D.~Stosic, N.~Burgess, M.~Cornea, P.~Dubey,
      R.~Grisenthwaite, S.~Ha, A.~Heinecke, P.~Judd, J.~Kamalu, et~al.
      \emph{FP8 Formats for Deep Learning}. arXiv:2209.05433, 2022.

    \bibitem{ocp2023fp8}
      Open Compute Project.
      \emph{OCP 8-bit Floating Point Specification (OFP8), Revision 1.0}, 2023.

    \bibitem{deepseek2024v3}
      DeepSeek-AI.
      \emph{DeepSeek-V3 Technical Report}. arXiv:2412.19437, 2024.

    \bibitem{deepgemm}
      DeepSeek-AI.
      \emph{DeepGEMM}. Repositorio de código,
      \url{https://github.com/deepseek-ai/DeepGEMM}.

    \bibitem{deepseek2025v31}
      DeepSeek-AI.
      \emph{DeepSeek-V3.1} (ficha del modelo). Hugging Face, 2025.
      \url{https://huggingface.co/deepseek-ai/DeepSeek-V3.1}.

  \end{thebibliography}

\end{document}
